\section*{Chapitre \ref{ch:b1:first-law} --- Le premier principe de la thermodynamique}

\begin{solution}{exo:b1:first-law:1}
$Q = C_{P,m}\Delta T = 29.1 \times 100 = \qty{2.91}{kJ}$~; $W = -nR\Delta T =
\qty{-0.83}{kJ}$~; $\Delta U = Q + W = \qty{2.08}{kJ} = C_{V,m}\Delta T$.
\end{solution}

\begin{solution}{exo:b1:first-law:2}
$W = -nRT\ln(V_2/V_1) = -2494\ln(0.25) = \qty{+3.46}{kJ}$~; $\Delta U = 0$ donc
$Q = \qty{-3.46}{kJ}$ (cédée au \omterm{def:b1:first-law:heat}{thermostat}).
\end{solution}

\begin{solution}{exo:b1:first-law:3}
$T_2 = T_1 \times 10^{\gamma - 1} = 300 \times 10^{0.4} = \qty{754}{K}$~; $P_2 =
10^{1.4} = \qty{25}{bar}$.
\end{solution}

\begin{solution}{exo:b1:first-law:4}
Argon~: $C_{V,m} = \tfrac32 R = 12.5$, $C_{P,m} = \tfrac52 R = 20.8$ J/(mol K),
$\gamma = 1.67$~; azote~: $20.8$, $29.1$, $\gamma = 1.40$. Pour l'eau, la
dilatation au chauffage est minime, donc le travail $P\,\dd V$ à pression
constante est négligeable devant la \omterm{def:b1:first-law:heat}{chaleur}~: $c_P - c_V \approx 0$.
\end{solution}

\begin{solution}{exo:b1:first-law:5}
$T_f = (200 \times 80 + 300 \times 20)/500 = \qty{44}{\degreeCelsius}$. Avec le
récipient~: $(200 \times 4.18 \times 80 + (300 \times 4.18 + 80) \times 20)/(500 \times
4.18 + 80) = \qty{43}{\degreeCelsius}$.
\end{solution}

\begin{solution}{exo:b1:first-law:6}
Glace~: réchauffement jusqu'à $0$~: $0.050 \times 2100 \times 10 = \qty{1.05}{kJ}$~;
fusion~: \qty{16.7}{kJ}~; au total \qty{17.8}{kJ}. L'eau peut céder \qty{25.1}{kJ}
en descendant à $0$~: tout fond, puis $0.250 \times 4180\,T_f = 25.1 - 17.8$ kJ~: $T_f =
\qty{7}{\degreeCelsius}$. Avec \qty{150}{g}~: il faudrait \qty{53.3}{kJ} $>$ \qty{25.1}{kJ}~:
seuls $(25.1 - 3.15)/334 = \qty{66}{g}$ de glace fondent, $T_f = \qty{0}{\degreeCelsius}$,
et \qty{84}{g} de glace subsistent.
\end{solution}

\begin{solution}{exo:b1:first-law:7}
$W = -P_1(V_2 - V_1) + 0 - P_2(V_1 - V_2) + 0 = (P_2 - P_1)(V_2 - V_1) > 0$~:
le cycle est parcouru dans le sens trigonométrique (détente à basse
pression, compression à haute), donc le gaz \emph{reçoit} un travail net,
égal à l'aire enclose~; dans l'autre sens il le fournirait --- un moteur.
\end{solution}

\begin{solution}{exo:b1:first-law:8}
$T_2 = 300 \times 7^{0.286} = \qty{523}{K}$. L'air comprimé et chaud échauffe
le corps de la pompe par conduction à chaque coup. Dans le pneu, l'air
refroidit jusqu'à l'ambiante à volume fixé~: $P \propto T$, la pression
baisse dans le rapport des températures --- il faut regonfler.
\end{solution}

\begin{solution}{exo:b1:first-law:9}
$W = 0$ (aucune pression extérieure à repousser), $Q = 0$~: $\Delta U = 0$,
donc $\Delta T = 0$ pour un \omterm{def:b1:kinetic-theory:model}{gaz parfait}. \omterm{rem:b1:first-law:irreversible}{Transformation irréversible}~: le gaz
ne revient jamais de lui-même. Un gaz réel, dont les molécules s'attirent,
se refroidit légèrement (une part de l'\omterm{thm:b1:work-and-energy:ke}{énergie cinétique} devient potentielle
lorsqu'elles s'éloignent).
\end{solution}

\begin{solution}{exo:b1:first-law:10}
$Q = L_v = \qty{2.26}{MJ}$~; $W = -P(V_g - V_l) = -1.013\times10^5 \times 1.67 =
\qty{-0.17}{MJ}$~; $\Delta U = \qty{2.09}{MJ}$~: $93\%$ de la \omterm{def:b1:first-law:heat}{chaleur} sert à
séparer les molécules (\omterm{def:b1:kinetic-theory:U}{énergie interne}), $7\%$ à repousser l'atmosphère.
\end{solution}

\begin{solution}{exo:b1:first-law:11}
$W = -P_2(V_2 - V_1)$, $\Delta U = nC_{V,m}(T_2 - T_1)$, $V_2 = nRT_2/P_2$, $V_1
= nRT_1/P_1$~: $\tfrac52(T_2 - 300) = -T_2 + 300 \times 5$, $3.5T_2 = 2250$,
$T_2 = \qty{643}{K}$, $V_2 = \qty{10.7}{L}$. Cas réversible~: $T_2 = 300 \times
5^{0.286} = \qty{475}{K}$, $V_2 = \qty{7.9}{L}$. La compression brusque
fournit plus de travail au gaz (les \qty{5}{bar} agissent dès le départ)~: il
finit plus chaud et moins comprimé.
\end{solution}

\begin{solution}{exo:b1:first-law:12}
$\rho = PM/RT = \qty{1.20}{kg/m^3}$~: $\sqrt{\gamma P/\rho} = \sqrt{1.4 \times
1.013\times10^5/1.20} = \qty{343}{m/s}$~; $\sqrt{P/\rho} = \qty{290}{m/s}$, soit $15\%$
trop bas. Avec $P/\rho = RT/M$~: $c = \sqrt{\gamma RT/M}$~; face à $u = \sqrt{3RT/M}$
= \qty{502}{m/s}~: $c/u = \sqrt{\gamma/3} = 0.68$.
\end{solution}

\begin{solution}{pb:b1:first-law:1}
\textbf{1.} $n = P_1V_1/RT_1 = 10^5 \times 2.12\times10^{-4}/2494 = \qty{8.5e-3}{mol}$.

\textbf{2.} $PV^\gamma$ constant~: $V_2 = V_1(P_1/P_2)^{1/\gamma} = 212 \times
7^{-0.714} = \qty{53}{cm^3}$.

\textbf{3.} $T_2 = T_1(P_2/P_1)^{(\gamma - 1)/\gamma} = 300 \times 7^{0.286} =
\qty{523}{K}$.

\textbf{4.} $W = nC_{V,m}(T_2 - T_1) = 8.5\times10^{-3} \times 20.8 \times 223 =
\qty{39}{J}$.

\textbf{5.} Air à ajouter~: $\Delta n = \Delta P\,V/RT = 6\times10^5 \times 2\times10^{-3}/
2494 = \qty{0.48}{mol}$~: $0.48/8.5\times10^{-3} = 57$ coups de pompe, environ
\qty{2.2}{kJ} --- une minute d'effort honnête.

\textbf{6.} À volume constant, $P \propto T$~: la pression baisse à mesure que
l'air refroidit (jusqu'au facteur $300/523$ pour l'air du dernier coup)~; le
cycliste redonne quelques coups après une pause.

\textbf{7.} $W_{\text{iso}} = nRT\ln(P_2/P_1) = 8.5\times10^{-3} \times 2494 \times
1.95 = \qty{41}{J} > \qty{39}{J}$~: le chemin isotherme comprime davantage le
gaz (jusqu'à \qty{30}{cm^3} au lieu de \qty{53}{cm^3}) pour atteindre \qty{7}{bar},
parce qu'il reste froid --- plus de volume balayé, donc plus de travail.

\textbf{8.} $n = 10^5 \times 5\times10^{-4}/2494 = \qty{0.020}{mol}$~; $T_2 =
300 \times 20^{0.4} = \qty{994}{K}$.

\textbf{9.} $P_2 = 20^{1.4} = \qty{66}{bar}$.

\textbf{10.} \qty{994}{K}, c'est bien au-dessus du point d'auto-inflammation
du carburant~: celui-ci s'enflamme dès son injection. Un moteur à essence ne
comprime son mélange air--carburant que jusqu'à \qty{750}{K} environ, pour
qu'il ne s'enflamme pas avant l'étincelle~; un taux trop élevé donne le «~cliquetis~».

\textbf{11.} $W = nC_{V,m}(T_2 - T_1) = 0.020 \times 20.8 \times 694 = \qty{289}{J}$.

\textbf{12.} $25 \times 289 = \qty{7.2}{kW}$ par cylindre, stockés dans le gaz
chaud et restitués à la détente.

\textbf{13.} L'irréversibilité (le gaz ne suit pas le chemin quasi statique~;
un travail supplémentaire est dissipé) élève $T_2$~; la \omterm{def:b1:first-law:heat}{chaleur} perdue aux
parois l'abaisse. Dans un moteur réel, le second effet l'emporte en général~:
$T_2$ est un peu au-dessous de la valeur idéale.

\textbf{14.} $Q = 20\times10^{-6} \times 43\times10^6 = \qty{860}{J}$.

\textbf{15.} Isobare~: $Q = nC_{P,m}(T_3 - T_2)$~: $T_3 - T_2 = 860/(0.020 \times
29.1) = \qty{1480}{K}$, $T_3 = \qty{2470}{K}$~; $V_3 = V_2T_3/T_2 = 25 \times
2.48 = \qty{62}{cm^3}$.

\textbf{16.} $W = -P_2(V_3 - V_2) = -66\times10^5 \times 37\times10^{-6} =
\qty{-244}{J}$~: travail fourni par le gaz.

\textbf{17.} $T_4 = T_3(V_3/V_1)^{\gamma - 1} = 2470 \times (0.124)^{0.4} =
\qty{1070}{K}$~; $W = nC_{V,m}(T_4 - T_3) = 0.020 \times 20.8 \times (-1400) =
\qty{-582}{J}$.

\textbf{18.} Travail net fourni $= 244 + 582 - 289 = \qty{537}{J}$~; rapport
$537/860 = 62\%$ (le rendement idéal du cycle Diesel avec ces réglages~; les
moteurs réels atteignent environ $40\%$).

\textbf{19.} $1 - 300/2470 = 88\%$~: le cycle Diesel idéal reste nettement
au-dessous de la borne de Carnot, parce qu'il reçoit sa \omterm{def:b1:first-law:heat}{chaleur} sur toute une
plage de températures et non à la plus haute.

\textbf{20.} Sur un cycle, $\Delta U = 0$~: $Q_{\text{out}} = -(860 - 537) =
\qty{-323}{J}$, c'est-à-dire \qty{323}{J} qui partent à l'échappement~; vérification~:
$nC_{V,m}(T_4 - T_1) = 0.020 \times 20.8 \times 770 = \qty{320}{J}$.

\textbf{21.} L'oscillation est rapide (une seconde) devant la conduction
thermique à travers \qty{10}{L} d'air~: elle est adiabatique. En linéarisant
$PV^\gamma$~: $\dd P/P = -\gamma\,\dd V/V = -\gamma Sx/V$.

\textbf{22.} Force sur la bille $S\,\dd P = -\gamma PS^2x/V$~: $m\ddot x =
-(\gamma PS^2/V)x$, \omterm{prop:b1:point-kinematics:harmonic}{mouvement harmonique}, $\omega_0^2 = \gamma PS^2/mV$.

\textbf{23.} $\omega_0^2 = 1.4 \times 10^5 \times 4\times10^{-8}/(0.0165 \times 0.010)
= \qty{33.9}{s^{-2}}$~: $T = 2\pi/5.83 = \qty{1.08}{s}$. À partir de $T = \qty{1.10}{s}$~:
$\gamma = 4\pi^2mV/(T^2PS^2) = 4\pi^2 \times 0.0165 \times 0.010/(1.21 \times 10^5
\times 4\times10^{-8}) = 1.35$.

\textbf{24.} Le frottement amortit le mouvement et, s'il est sec, déplace la
position d'équilibre~; les fuites autour de la bille laissent l'air s'échapper
pendant une compression, ce qui affaiblit la force de rappel~: les deux
allongent la période et biaisent $\gamma$ vers le bas --- comme le $1.35$ ci-dessus.

\textbf{25.} Pompe et Diesel~: $Q = 0$, donc $\Delta U = W$ --- le travail est
devenu \omterm{def:b1:kinetic-theory:U}{énergie interne}, donc température, par la \omterm{prop:b1:first-law:transformations}{loi de Laplace} d'exposant
$\gamma$. Bille~: une compression rapide est adiabatique, et $\gamma$ fixe la
raideur d'un ressort d'air.
\end{solution}
